\documentclass[12pt]{article}
\usepackage[margin=1in]{geometry}
\usepackage{amsmath,amssymb,amsthm}
\usepackage{enumitem}

\newcommand{\N}{\mathbb{N}}
\newcommand{\R}{\mathbb{R}}
\newcommand{\Z}{\mathbb{Z}}

% Answer environment -- fill in your proofs here
\newenvironment{answer}{\begin{proof}[Answer]}{\end{proof}}

\title{Mathematics 629 --- Homework 7}
\author{Alexander Kim}
\date{Due: Monday, April 13, 2026}

\begin{document}
\maketitle

%%% Problem 1 %%%%%%%%%%%%%%%%%%%%%%%%%%%%%%%%%%%%%%%%%%%
\section*{Problem 1}
Let $X_1 = [0,1]$, with $\mathcal{M}_1 = \mathcal{B}_{[0,1]}$ and with $\mu_1$ the Lebesgue measure, $X_2 = [0,1]$ with $\mathcal{M}_2 = \mathcal{B}_{[0,1]}$ and $\mu_2$ being counting measure. Let $\mu = \mu_1 \times \mu_2$. Consider the diagonal $D = \{(x,x) : 0 \leq x \leq 1\}$.

\begin{enumerate}[label=(\roman*)]
	\item Use the definition of $\mu$ to evaluate $\int \mathbf{1}_D\,d\mu$.

	      \begin{answer}
		      By definition of product measure, $\int \mathbf{1}_D d\mu = \int \mu_2(sl_{x_1}) d\mu_1(x_1)$ where $sl_{x_1} = x_1$. Then $\int \mathbf{1}_D d\mu = \int_{[0,1]}1 d\mu_1 = 1.$
	      \end{answer}

	\item Show that $\int\!\int \mathbf{1}_D\,d\mu_1\,d\mu_2$ and $\int\!\int \mathbf{1}_D\,d\mu_2\,d\mu_1$ have different values. Why can Tonelli's theorem not be applied?

	      \begin{answer}
		      In (i), integrating with respect to $\mu_2$ first showed $\int \int \mathbf{1}_D d\mu_2 d\mu_1 = 1$. Let $sl_{x_2} = x_2$. Then similarly,
		      $$ \int \int \mathbf{1}_D d\mu_1 d\mu_2 = \int \mu_1(sl_{x_2}) d\mu_2(x_2) = 0*#(sl_{x_2}) = 0.$$
		      Clearly order of integration matters in this case. Since Tonelli guarantees equivalence regardless of order of integration, the theorem cannot be applied here because $(X_2, M_2, \mu_2)$ is not sigma finite $(\mu_2([0,1]$ is uncountable).
	      \end{answer}
\end{enumerate}

%%% Problem 2 %%%%%%%%%%%%%%%%%%%%%%%%%%%%%%%%%%%%%%%%%%%
\section*{Problem 2}
Consider, for $x \in (0,1) \times (0,1)$,
\[
	f(x_1, x_2) = \frac{x_1^2 - x_2^2}{(x_1^2 + x_2^2)^2} = \frac{\partial^2}{\partial x_1\,\partial x_2}\arctan\frac{x_1}{x_2}.
\]

\begin{enumerate}[label=(\roman*)]
	\item Verify that
	      \[
		      \int_0^1\!\int_0^1 \frac{x_1^2 - x_2^2}{(x_1^2 + x_2^2)^2}\,dx_2\,dx_1 = \frac{\pi}{4}, \qquad
		      \int_0^1\!\int_0^1 \frac{x_1^2 - x_2^2}{(x_1^2 + x_2^2)^2}\,dx_1\,dx_2 = -\frac{\pi}{4}.
	      \]

	      \begin{answer}
		      \begin{enumerate}
			      \item First observe that since $\frac{\partial^2}{\partial x_1 \partial x_2}\arctan \frac{x_1}{x_2} = f$ then $$\frac{\partial}{\partial x_2} \arctan \frac{x_1}{x_2} = \frac{1}{1 + (x_1/x_2^2)} * (\frac{-x_1}{x_2^2}) = \frac{-x_1}{x_1^2+x_2^2}.$$
			            In other words,
			            $$\int_0^1 \int_0^1 f dx_2 dx_1 = \int_0^1 \frac{1}{1+x_1^2} dx_1 = \arctan(x_1)|_0^1 = \frac{\pi}{4}.$$

			            Similarly, integrating in the other order shows
			            $$\frac{\partial}{\partial x_1} \frac{-x_1}{x_1^2 + x_2^2} = f$$ so then
			            by keeping $x_1$ fixed we get that
			            $$\int_0^1 \int_0^1 f dx_1 dx_2 = \int_0^1 \frac{-1}{1+ x_2^2} dx_2 = -\arctan(x_2)|_0^1\frac{-\pi}{4}.$$
		      \end{enumerate}
	      \end{answer}

	\item Let $T = \{x : 0 < x_2 < x_1 < 1\}$, let $m$ be Lebesgue measure in $(0,1)^2$ and show that
	      \[
		      \int_T f\,dm = \infty.
	      \]

	      \begin{answer}
		      Observe that using $\frac{\partial}{\partial x_2}(\frac{x_2}{x_1^2+x_2^2} = f$ then
		      $$\int_0^{x_1} = [\frac{x_2}{x_1^2+x_2^2}]_{x_2 = 0}^{x_2=x_1} = \frac{1}{2x_1}.$$
		      Then integrating over $x_1$ shows
		      $$\int_T f dm = \int_0^1 \frac{1}{2x_1}dx_1 = \frac{1}{2}[\ln x_1]_0^1 = \infty$$.
	      \end{answer}

	\item Why can Fubini's theorem not be applied in (i)?

	      \begin{answer}
		      Fubinis theorem requires that $f\in L^1((0,1)^2)$ but we have shown this isn't true in part (ii). This is why part (i) yields integrals of different values without contradicting Fubini's Theorem.
	      \end{answer}
\end{enumerate}

%%% Problem 3 %%%%%%%%%%%%%%%%%%%%%%%%%%%%%%%%%%%%%%%%%%%
\section*{Problem 3}
Let $(X_1, \mathcal{M}_1, \mu_1)$ be a measure space and $X_2 = \N$ with $\mathcal{M}_2 = \mathcal{P}(\N)$ and counting measure $\mu_2$. Interpret Tonelli's and Fubini's theorem in this situation and show that one does not need the assumption of $\sigma$-finiteness of $\mu_1$ in this case.

\begin{answer}
	Let $f: X\times \mathbb{N} \to \mathbb{R}$ and let $f(n(x_1) = f(x_1, n)$ for $x_1\in X_1$ and $n\in \mathbb{N}$. Recall that counting measure is a sum so
	$$\int_{\mathbb{N}}f(x_1,n)d\mu_2 = \sum_{n=1}^\infty f(n(x_1).$$
	By Tonelli's theorem,
	$$\int_{X_1} \sum_{n=1}^\infty f_n(x_1)d\mu_1 = \sum_{n=1}^\infty \int_{X_1} f(n)x_1)d\mu_1.$$ However this result follows as a result of MCT. If we let $S_n = \sum_{n=1}^N f(n(x_1)$ be the partial sum up to some $N$ then by MCT
	$$\int_{X_1} \sum_{n=1}^\infty f_n d\mu_1 = lim_{N\to \infty} \int_{X_1} \sum_{n=1}^N f_n d\mu_1 = lim_{N\to \infty}\sum_{n=1}^N \int_{X_1}f_n d\mu_1 = \sum_{n=1}^\infty \int_{X_1} f_n d\mu_1.$$
	This proof doesn't rely on the $\sigma$-finiteness of $\mu_1$ and furthermore, to show $\sum_N \int |f_n|d\mu_1 < \infty$ apply Tonelli's theorem to $f_n^+$ and $f_n^-$. In both cases, $\sigma$-finiteness of $\mu_1$ is not needed since sums over $\mathbb{N}$ are always well defined.
\end{answer}

%%% Problem 4 %%%%%%%%%%%%%%%%%%%%%%%%%%%%%%%%%%%%%%%%%%%
\section*{Problem 4}
Let $(X, \mathcal{M}, \mu)$ be a $\sigma$-finite measure space. Let $f$ be a nonnegative measurable function on $X$ and
\[
	U_f = \{(x,y) \in X \times [0,\infty] : 0 \leq y \leq f(x)\}.
\]
Show that $U_f$ is $\mathcal{M} \times \mathcal{B}_{\R}$ measurable and $\mu \times m(U_f) = \int f\,d\mu$. Show that the Lebesgue measure (in $\R^2$) of the graph of $f$ (i.e., $\{(x, f(x)) : x \in X\}$) is zero.

\begin{answer}
	Since we have $\sigma$-finiteness then let $sl_x = sl_x([0,\infty]) =\{y: (x, y) \in U_f \}$. Then integrating over $sl_x$ observe that
	$$\int_0^\infty \mathbf{1}_{U_f}(sl_x)dm(y) = f(x)$$
	since $$\mu \times m(U_f) = \int_0^\infty \int_X \mathbf{1}_{U_f} d(\mu \times m).$$
	In other words, $\mu \times m(U_f) = \int_X f(x) d\mu(x) = \int f d\mu.$
	Let $G_f = \{(x, f(x)): x \in X\}$ be the graph of $f$. Note that $sl_y(G_f) = \{f(x)\}$ which as a singleton has $m(sl_y) = 0$. Thus
	$$\mu \times m(G_f) = \int_X m(sl_y)d\mu(x) = \int_X 0 d\mu = 0.$$
	So the measure of the graph of $f$ is 0.
\end{answer}

%%% Problem 5 %%%%%%%%%%%%%%%%%%%%%%%%%%%%%%%%%%%%%%%%%%%
\section*{Problem 5}
Let
\[
	A_1 := \{(x,y,z) \in \R^3 : x^2 + z^2 \leq 1\}, \qquad A_2 := \{(x,y,z) \in \R^3 : y^2 + z^2 \leq 1\}
\]
and compute the Lebesgue measure of $A_1 \cap A_2$.

\begin{answer}
	Define $sl_{xy}(E) \in \mathbb{R^2}$ as the generalized rectangle for fixed $z$ given some $E \in A_1 \cap A_2$ and $sl_x(E) \in \mathbb{R}$ as the generalized rectangle for fixed $y$ and $z$. Note that $m = m_1 \times m_2 \times m_3$ where $m_1, m_2, m_3$ are Lebesgue measure. Since Lebesgue measure is $\sigma$-finite, by the generalized Cavelieri principle we can compute
	\begin{align}
		m(A_1 \cap A_2) & = \int m_1\times m_2(sl_{xy})dm_3(z)                                          \\
		                & = \int \int m_2(sl_x(E))dm_1(x)dm_3(z)                                        \\
		                & = \int_{-1}^1 \int_{\sqrt{1-z^2}}^{\sqrt{1-z^2}} 2\sqrt{1-z^2} dm_1(x)dm_3(z) \\
		                & = \int_{-1}^1 4(1-z^2)dm_3(z) = \frac{16}{3}.
	\end{align}
\end{answer}

\end{document}
